\chapter{Canonical Quantization of the Real Scalar Field}

\begin{remarkbox}
This chapter turns the oscillator picture of Chapter 1 into the standard
canonical quantization of the free real scalar field. The main new object is not
one oscillator coordinate, but a field operator \(\hat{\phi}(t,\vect{x})\)
with one canonical commutation relation at every point of space.
\end{remarkbox}

\section{Equal-Time Commutation Relations}

We work with a real scalar field in units \(c=\hbar=1\). With metric convention
\((+,-,-,-)\), the free Lagrangian density is
\[
    \mathcal{L}
    = \frac{1}{2}\partial_\mu\phi\,\partial^\mu\phi
      - \frac{1}{2}m^2\phi^2
    = \frac{1}{2}\dot\phi^2
      - \frac{1}{2}(\nabla\phi)^2
      - \frac{1}{2}m^2\phi^2.
\]
The canonical momentum conjugate to \(\phi\) is
\[
    \pi(t,\vect{x})
    = \frac{\partial\mathcal{L}}{\partial\dot\phi(t,\vect{x})}
    = \dot\phi(t,\vect{x}).
\]
Classically, the canonical Poisson bracket is
\[
    \{\phi(t,\vect{x}),\pi(t,\vect{y})\}_{\mathrm{PB}}
    = \delta^{(3)}(\vect{x}-\vect{y}).
\]
Canonical quantization replaces the fields by operators and the Poisson bracket
by a commutator:
\[
    [\hat{\phi}(t,\vect{x}),\hat{\pi}(t,\vect{y})]
    = i\delta^{(3)}(\vect{x}-\vect{y}).
\]
The remaining equal-time commutators are
\[
    [\hat{\phi}(t,\vect{x}),\hat{\phi}(t,\vect{y})]=0,
    \qquad
    [\hat{\pi}(t,\vect{x}),\hat{\pi}(t,\vect{y})]=0.
\]
These are the field-theoretic version of the oscillator relation
\([\hat q,\hat p]=i\). Instead of one coordinate and one momentum, there is one
canonical pair at every spatial point.

\begin{definitionbox}
The \emph{equal-time canonical commutation relations} are the commutators
between field and momentum operators evaluated at the same time. They are the
canonical starting point of the operator quantization of a bosonic field.
\end{definitionbox}

\section{Field Operator and Momentum Operator}

The field operator \(\hat{\phi}(t,\vect{x})\) is the quantum object corresponding
to the classical field \(\phi(t,\vect{x})\). Its conjugate momentum operator is
\[
    \hat{\pi}(t,\vect{x})=\partial_t\hat{\phi}(t,\vect{x}).
\]
The free equation of motion is the Klein--Gordon equation,
\[
    (\Box+m^2)\hat{\phi}(x)=0,
\]
where \(x=(t,\vect{x})\). This equation has the same form as the classical field
equation, but its meaning is different: it is an equation for an operator-valued
object.

The Hamiltonian density of the free field is
\[
    \mathcal{H}
    = \frac{1}{2}\pi^2
      + \frac{1}{2}(\nabla\phi)^2
      + \frac{1}{2}m^2\phi^2.
\]
After quantization one writes formally
\[
    \hat H
    = \int \dd^3x\,
    \left(
        \frac{1}{2}\hat{\pi}^2
        + \frac{1}{2}(\nabla\hat{\phi})^2
        + \frac{1}{2}m^2\hat{\phi}^{\,2}
    \right).
\]
The word ``formally'' is important. Products such as
\(\hat{\phi}(x)^2\) at the same point require care. This issue will return when
we discuss normal ordering and renormalization.

\section{Mode Expansion of the Field Operator}

The free Klein--Gordon equation admits plane-wave solutions. The relativistic
energy of a mode with spatial momentum \(\vect{k}\) is
\[
    \omega_{\vect{k}}=\sqrt{\vect{k}^{\,2}+m^2}.
\]
A convenient expansion of the real scalar field operator is
\[
    \hat{\phi}(t,\vect{x})
    = \int \frac{\dd^3k}{(2\pi)^3}
      \frac{1}{\sqrt{2\omega_{\vect{k}}}}
      \left(
          \hat a_{\vect{k}}
          e^{-i\omega_{\vect{k}}t+i\vect{k}\cdot\vect{x}}
          +
          \hat a_{\vect{k}}^\dagger
          e^{i\omega_{\vect{k}}t-i\vect{k}\cdot\vect{x}}
      \right).
\]
The corresponding momentum operator is obtained by differentiating with respect
to time:
\[
    \hat{\pi}(t,\vect{x})
    = -i\int \frac{\dd^3k}{(2\pi)^3}
      \sqrt{\frac{\omega_{\vect{k}}}{2}}
      \left(
          \hat a_{\vect{k}}
          e^{-i\omega_{\vect{k}}t+i\vect{k}\cdot\vect{x}}
          -
          \hat a_{\vect{k}}^\dagger
          e^{i\omega_{\vect{k}}t-i\vect{k}\cdot\vect{x}}
      \right).
\]
The normalization factors are chosen so that the canonical equal-time
commutation relations take their standard form.

Because the field is real, the operator \(\hat{\phi}\) is Hermitian:
\[
    \hat{\phi}(t,\vect{x})^\dagger=\hat{\phi}(t,\vect{x}).
\]
This is why the expansion contains both \(\hat a_{\vect{k}}\) and
\(\hat a_{\vect{k}}^\dagger\). The negative-frequency part is the Hermitian
conjugate of the positive-frequency part.

\section{Creation and Annihilation Operators for Momentum Modes}

The operators \(\hat a_{\vect{k}}\) and \(\hat a_{\vect{k}}^\dagger\) satisfy
\[
    [\hat a_{\vect{k}},\hat a_{\vect{p}}^\dagger]
    = (2\pi)^3\delta^{(3)}(\vect{k}-\vect{p}),
\]
with
\[
    [\hat a_{\vect{k}},\hat a_{\vect{p}}]=0,
    \qquad
    [\hat a_{\vect{k}}^\dagger,\hat a_{\vect{p}}^\dagger]=0.
\]
These relations are the continuum analogue of
\([\hat a_{\vect{k}},\hat a_{\vect{p}}^\dagger]
=\delta_{\vect{k},\vect{p}}\) in a finite box.

The interpretation is inherited from the harmonic oscillator. The operator
\(\hat a_{\vect{k}}^\dagger\) creates one quantum of momentum \(\vect{k}\), while
\(\hat a_{\vect{k}}\) annihilates one such quantum. The word ``quantum'' here
means an excitation of the free scalar field mode labelled by \(\vect{k}\).

\begin{examplebox}
If the field is placed in a finite box, the momentum labels are discrete and the
commutator is \([\hat a_{\vect{k}},\hat a_{\vect{p}}^\dagger]
=\delta_{\vect{k},\vect{p}}\). In infinite volume, the Kronecker delta is
replaced by \((2\pi)^3\delta^{(3)}(\vect{k}-\vect{p})\).
\end{examplebox}

\section{Fock Space}

The free-field vacuum is defined by
\[
    \hat a_{\vect{k}}|0\rangle=0
    \qquad \text{for every } \vect{k}.
\]
The Fock space is built by acting on this vacuum with creation operators. A
one-particle state is
\[
    |\vect{k}\rangle
    = \hat a_{\vect{k}}^\dagger|0\rangle.
\]
A two-particle state is
\[
    |\vect{k}_1,\vect{k}_2\rangle
    = \hat a_{\vect{k}_1}^\dagger
      \hat a_{\vect{k}_2}^\dagger|0\rangle,
\]
and similarly for higher particle numbers.

For a real scalar field the quanta are neutral bosons. The creation operators
commute with one another, so exchanging two creation operators does not change
the state. This is the operator origin of Bose symmetry in the free scalar
field.

\begin{definitionbox}
The \emph{Fock space} of the free real scalar field is the direct sum of sectors
with zero, one, two, and more scalar quanta. It is generated from the vacuum by
all finite products of creation operators.
\end{definitionbox}

\section{Vacuum, One-Particle States, and Many-Particle States}

The vacuum \(|0\rangle\) is not a classical field configuration. It is a vector
in Hilbert space. It is characterized by the absence of free quanta, not by the
absence of operators or fluctuations.

A one-particle state \(|\vect{k}\rangle\) is an eigenstate of momentum and
energy in the free theory. A general one-particle wave packet is a superposition
of momentum eigenstates,
\[
    |f\rangle
    = \int \frac{\dd^3k}{(2\pi)^3}\,f(\vect{k})
      \hat a_{\vect{k}}^\dagger|0\rangle.
\]
The function \(f(\vect{k})\) controls the momentum profile of the state. A
localized particle-like state is not a single momentum eigenstate, but a wave
packet built from many momenta.

Many-particle states are obtained by applying several creation operators. For
example,
\[
    |f,g\rangle
    = \int \frac{\dd^3k}{(2\pi)^3}
      \int \frac{\dd^3p}{(2\pi)^3}
      f(\vect{k})g(\vect{p})
      \hat a_{\vect{k}}^\dagger\hat a_{\vect{p}}^\dagger|0\rangle.
\]
Because the field is bosonic, the state is symmetric under exchange of the two
created quanta.

\section{Hamiltonian and Momentum Operators}

Substituting the mode expansion into the free Hamiltonian gives
\[
    \hat H
    = \int \frac{\dd^3k}{(2\pi)^3}\,
      \omega_{\vect{k}}
      \left(
          \hat a_{\vect{k}}^\dagger\hat a_{\vect{k}}
          + \frac{1}{2}(2\pi)^3\delta^{(3)}(0)
      \right).
\]
The term proportional to \(\delta^{(3)}(0)\) is the infinite-volume version of
the zero-point energy. In a finite box, the same expression appears as
\[
    E_0=\frac{1}{2}\sum_{\vect{k}}\omega_{\vect{k}}.
\]

The spatial momentum operator is
\[
    \hat{\vect{P}}
    = \int \dd^3x\,\hat{\pi}\nabla\hat{\phi}.
\]
In terms of the ladder operators it becomes, after the usual ordering convention,
\[
    \hat{\vect{P}}
    = \int \frac{\dd^3k}{(2\pi)^3}\,
      \vect{k}\,\hat a_{\vect{k}}^\dagger\hat a_{\vect{k}}.
\]
Thus the state \(\hat a_{\vect{k}}^\dagger|0\rangle\) has momentum \(\vect{k}\)
and energy \(\omega_{\vect{k}}\), exactly as required by the relativistic
relation
\[
    \omega_{\vect{k}}^2=\vect{k}^{\,2}+m^2.
\]

\section{Normal Ordering}

The infinite constant in the Hamiltonian motivates a useful convention. The
normal-ordered form of an operator is obtained by moving all creation operators
to the left of all annihilation operators. For the oscillator product one writes
\[
    :\hat a_{\vect{k}}\hat a_{\vect{p}}^\dagger:
    = \hat a_{\vect{p}}^\dagger\hat a_{\vect{k}}.
\]
The normal-ordered Hamiltonian of the free scalar field is
\[
    :\hat H:
    = \int \frac{\dd^3k}{(2\pi)^3}\,
      \omega_{\vect{k}}
      \hat a_{\vect{k}}^\dagger\hat a_{\vect{k}}.
\]
It satisfies
\[
    \langle 0|:\hat H:|0\rangle=0.
\]

Normal ordering is not a universal cure for all infinities. It is a clean way
to remove the constant vacuum energy of a free field when only energy
differences above the free vacuum are relevant. Interacting theories require a
more systematic renormalization procedure.

\section{Zero-Point Energy and Its Interpretation}

Every harmonic oscillator contributes \(\omega/2\) to its ground-state energy.
A field has one oscillator for every momentum mode, so the formal zero-point
energy is
\[
    E_0=\frac{1}{2}\sum_{\vect{k}}\omega_{\vect{k}}
\]
in a finite box, or
\[
    E_0=\frac{V}{2}\int\frac{\dd^3k}{(2\pi)^3}\,\omega_{\vect{k}}
\]
in infinite-volume notation. This expression diverges at large
\(|\vect{k}|\).

In non-gravitational scattering calculations, a constant shift of the energy is
usually unobservable, so normal ordering is often sufficient at the free-field
level. Conceptually, however, the divergence is important. It shows that a field
has infinitely many short-distance degrees of freedom. Later chapters will treat
such divergences by regularization and renormalization rather than by simply
ignoring them.

\section{Field Operators as Operator-Valued Distributions}

The notation \(\hat{\phi}(x)\) suggests an operator at a spacetime point. This
is useful notation, but it hides a subtle fact: quantum fields are really
operator-valued distributions. They become well-defined operators only after
being smeared with suitable test functions.

Given a smooth function \(f(\vect{x})\), define a smeared field at fixed time by
\[
    \hat{\phi}(f)
    = \int \dd^3x\, f(\vect{x})\hat{\phi}(t,\vect{x}).
\]
This object is better behaved than the formal point operator
\(\hat{\phi}(t,\vect{x})\). Products of fields at the same point, such as
\(\hat{\phi}(x)^2\), are generally singular and require definition by a
limiting or renormalization procedure.

This is not merely a mathematical technicality. The singular behavior of local
products is one of the reasons renormalization is unavoidable in interacting
quantum field theory.

\section{Causality and Spacelike Commutators}

Relativistic QFT must respect the causal structure of spacetime. For a scalar
field, the basic causality condition is that field operators commute at
spacelike separation:
\[
    [\hat{\phi}(x),\hat{\phi}(y)]=0
    \qquad \text{if}\qquad (x-y)^2<0.
\]
This property is called microcausality. It means that measurements associated
with spacelike separated regions cannot influence one another through the order
in which the corresponding operators are applied.

The equal-time commutator
\[
    [\hat{\phi}(t,\vect{x}),\hat{\phi}(t,\vect{y})]=0
\]
is the simplest special case. The nontrivial statement is that the commutator
also vanishes for all spacelike separations, not only for equal times.

For the free scalar field, the commutator can be written as a c-number function,
\[
    [\hat{\phi}(x),\hat{\phi}(y)] = i\Delta(x-y),
\]
where \(\Delta\) is the Pauli--Jordan commutator function. Lorentz invariance
and the mass-shell condition imply that \(\Delta(x-y)=0\) outside the light cone.
This is the operator statement that the quantized scalar field is compatible
with relativistic causality.

\begin{summarybox}
Canonical quantization of the real scalar field begins with equal-time
commutators for \(\hat{\phi}\) and \(\hat{\pi}\). The free field is expanded in
momentum modes with creation and annihilation operators. Acting with creation
operators on the vacuum builds Fock space, while the Hamiltonian and momentum
operators assign energy \(\omega_{\vect{k}}\) and momentum \(\vect{k}\) to each
one-particle state. Normal ordering removes the free vacuum constant, but the
deeper lesson is that local quantum fields are operator-valued distributions and
must satisfy microcausality at spacelike separation.
\end{summarybox}

\section{Chapter Checklist}

After reading this chapter, one should be able to:
\begin{itemize}
    \item write the equal-time commutation relations of the real scalar field;
    \item derive the canonical momentum from the free scalar Lagrangian;
    \item explain the mode expansion of \(\hat{\phi}\) and \(\hat{\pi}\);
    \item state the commutators of \(\hat a_{\vect{k}}\) and
    \(\hat a_{\vect{k}}^\dagger\);
    \item construct vacuum, one-particle, and many-particle states;
    \item identify the Hamiltonian and momentum operators in terms of number
    operators;
    \item explain what normal ordering does and what it does not do;
    \item explain why local fields should be treated as distributions;
    \item state the microcausality condition for a relativistic scalar field.
\end{itemize}
